\subsection{Bonds with embedded options}

These entries value callable, putable, capped and floored bonds on the calibrated tree of \hyperref[meth:binomial-rate-tree-calibration]{the rate tree}, and measure their spread and their interest-rate risk. Convertible bonds, compared with straight bonds and stock, are valued in \hyperref[meth:convertible-bond-valuation]{convertible bonds}, in the options chapter. The examples use the Handbook's tree (par rates 3.5 to 5.2 percent, 10 percent volatility) so that each figure can be checked against the source; numbers are verified in \texttt{trading\_notes.numerical.embedded\_option\_bonds.rate\_tree\_checks}.

\begin{qfmethod}[Callable and putable bonds: decomposition into a straight bond and an option, valuation on a rate tree, and the effects of volatility and of the level and shape of the curve]\label{meth:callable-putable-bond-valuation}
\qffield{Statement} \emph{Securities with embedded options.} Callable bonds let the issuer redeem early, putable bonds let the holder sell back, extendible bonds let the holder extend maturity (economically a putable bond), sinking-fund bonds retire principal on a schedule, capped and floored floaters limit their coupons, and convertibles add an option on the issuer's shares (\hyperref[def:bond-cash-flow-contingencies]{cash-flow contingencies}; \hyperref[meth:capped-floored-floater-valuation]{floaters}; \hyperref[meth:convertible-bond-valuation]{convertibles}). \emph{Decomposition.} A callable bond is a straight bond less the call option the investor has sold to the issuer; a putable bond is a straight bond plus the put the investor holds (\hyperref[def:bond-cash-flow-contingencies]{cash-flow contingencies}). So $V_{\text{callable}}=V_{\text{straight}}-V_{\text{call}}$ and $V_{\text{putable}}=V_{\text{straight}}+V_{\text{put}}$.

\emph{Valuation on the tree.} Backward induction proceeds as for an option-free bond, but at every node where the option can be exercised the value is replaced by the exercise outcome: $V_t=\min[\text{call price},PV(\text{future cash flows})]$ for a call and $V_t=\max[\text{put price},PV]$ for a put. The option is worth the difference from the straight bond on the same tree \citep[chap.~36, pp.~840--845]{s7-fabozzifixed}. A bond that is both callable and putable tests the call first and then the put at each node (p.~845). A step-up callable note changes the coupon at each node according to its schedule (pp.~845--847).

\emph{Volatility.} Higher volatility raises the value of any option, so it lowers the callable's value and raises the putable's; the straight bond does not change. \emph{Level and shape of the curve.} Lower rates move the call into the money and the put out of it. A flatter curve whose later forward rates are lower makes the call more valuable, because exercise is likelier on the later nodes. Interpretation, checked numerically below.

\qffield{Practice} \emph{Reproduced} (Handbook, 6.5 percent four-year bond, exercisable at par after year 1, $\sigma=10$ percent).
\begin{itemize}
\item The straight bond is worth 104.643, the callable 102.899 and the putable 105.327. The call is worth 1.745 unrounded; the Handbook prints 1.744, from rounded values, and writes the difference as $104.634-102.899$, a transposition of 104.643. The put is worth 0.684.
\item At node HLL in year 3 the callable's continuation value is 100.315, so it is called and 100 enters the recursion (p.~841).
\item A step-up note paying 4.25 percent for two years and 7.50 percent for two is worth 102.082 noncallable and 100.031 callable at par at years 2 and 3, an option worth 2.051.
\end{itemize}

\emph{Checked effects} (original, same curve).
\begin{itemize}
\item \emph{Volatility.} Recalibrated at 20 percent volatility, the straight bond is still 104.643, but the callable falls to 102.108 and the putable rises to 106.010.
\item \emph{Level.} With par rates 100 basis points higher, the call is worth 0.419 instead of 1.745 and the put 1.669 instead of 0.684.
\item \emph{Shape.} On a flat 4.5 percent curve the call is worth 5.261, because the later rates the option depends on are lower than on the upward-sloping curve.
\end{itemize}

\qffield{Conditions} The exercise rule (call whenever the continuation value exceeds the call price) ignores notice periods, refunding costs and the issuer's own credit. The tree assumes one factor and a constant volatility; option values are model outputs and change with both.

\qffield{Pitfalls} Valuing a callable bond to its first call date or to maturity (yield to worst) and treating that as its value. Forgetting that the investor is short the call. Expecting a higher volatility to change the straight bond's value.

\qffield{See also} \hyperref[meth:option-adjusted-spread]{option-adjusted spread}; \hyperref[meth:embedded-option-effective-duration]{effective duration}; \hyperref[meth:capped-floored-floater-valuation]{capped and floored floaters}.

\end{qfmethod}

\begin{qfmethod}[Option-adjusted spread: the constant spread to the tree that reprices a bond with embedded options, how it is used, and why it falls as assumed volatility rises]\label{meth:option-adjusted-spread}
\qffield{Statement} The OAS is ``the fixed spread (usually measured in basis points) over the benchmark rates that equates the output from the valuation process with the actual market price of the security''. It is found by adding a constant spread to every rate on the calibrated tree and searching until the model value equals the price. Because the option's value depends on volatility, ``the OAS \dots\ is volatility-dependent'' \citep[chap.~36, pp.~849--850]{s7-fabozzifixed}. The spread can change exercise decisions on the tree, which is why the search needs a full revaluation at each trial (p.~849).

\emph{Use.} The OAS is the spread an investor earns for the bond's credit and liquidity risk after the option has been valued. It compares bonds with different options on a like-for-like basis. The \emph{Z-spread}, a constant spread over the spot curve that ignores the option, includes the option's cost, so for a callable bond option cost $\approx$ Z-spread $-$ OAS (Interpretation).

\emph{Volatility.} A higher assumed volatility raises the value of the call the investor has sold, lowers the model value of the callable and so leaves less spread to reprice it: the callable's OAS falls (Interpretation).

\qffield{Practice} \emph{Reproduced} (Handbook, pp.~849--850). The 6.5 percent callable bond is worth 102.899 on the tree but trades at 102.218. Adding 35 basis points to every rate reprices it; the solver gives 35.0 basis points. With the shift, node HLL is worth 99.985 and is no longer called (Exhibit~36-19).

\emph{Original.} The same callable has a Z-spread of 67.8 basis points, so the option costs about 32.8 basis points of spread. Assuming 20 percent volatility instead of 10, the OAS falls to $-6.0$ basis points: the bond looks rich rather than cheap. Two analysts who quote OAS at different volatilities are not comparing the same number.

\qffield{Conditions} The OAS depends on the benchmark curve, the volatility and the rate model; it is a relative-value measure under those assumptions, not a market observable. It assumes the spread is constant across nodes and dates.

\qffield{Pitfalls} Comparing OAS figures computed with different volatilities or benchmark curves. Reading a high Z-spread on a callable as cheapness when it pays for the call. Computing OAS without rerunning the exercise decisions at each trial spread.

\qffield{See also} \hyperref[def:bond-yield-spread-measures]{yield spreads}; \hyperref[meth:callable-putable-bond-valuation]{callable bonds}.

\end{qfmethod}

\begin{qfmethod}[Interest-rate risk of bonds with embedded options: effective duration and convexity at constant OAS, one-sided and key-rate durations, and callable, putable and straight bonds compared]\label{meth:embedded-option-effective-duration}
\qffield{Statement} \emph{Effective duration and convexity} for a bond with options use $D=(V_--V_+)/(2V_0\Delta r)$ and, in the convention used throughout this compendium (\hyperref[meth:convexity-portfolio-duration]{convexity}), $C=(V_-+V_+-2V_0)/(V_0\Delta r^2)$, with the price change $\Delta V/V_0\approx-D\,\Delta r+\tfrac12C\,\Delta r^2$; $V_\pm$ are values after the benchmark curve shifts by $\pm\Delta r$. The Handbook divides by $2V_0\Delta r^2$, which gives half the figure; with that convention the price change is $-D\,\Delta r+C_{\text{half}}\,\Delta r^2$. A convexity number must be quoted with its convention. The Handbook's procedure keeps the OAS constant (\hyperref[meth:option-adjusted-spread]{OAS}):
\begin{enumerate}
\item compute the OAS from the market price;
\item shift the par curve;
\item rebuild and recalibrate the tree;
\item add the OAS to every rate;
\item revalue the bond to get $V_+$ (and likewise $V_-$).
\end{enumerate}
The callable bond ``exhibits negative convexity'' \citep[chap.~36, pp.~850--851]{s7-fabozzifixed}. Yield-based measures do not capture options (\hyperref[meth:curve-based-empirical-duration]{curve-based measures}).

\emph{One-sided durations} use only the down-shift or only the up-shift: $D^-=(V_--V_0)/(V_0\Delta r)$ and $D^+=(V_0-V_+)/(V_0\Delta r)$. They differ when an option is near the money. A callable gains less when rates fall, because the call caps its price, than it loses when rates rise. A putable loses less when rates rise, because the put floors its price (Interpretation).

\emph{Key-rate durations} shift one par rate at a time, recalibrating the tree and holding the OAS. They show where on the curve the exposure sits.

\qffield{Practice} \emph{Reproduced} (Handbook, $\Delta r=25$ basis points, OAS 35 basis points): $V_+=101.621$, $V_-=102.765$, $V_0=102.218$, an effective duration of 2.24 and a negative convexity, printed as $-39.1$ in the Handbook's half-scaled convention from rounded prices. Unrounded it is $-38.574949$ in that convention and $-77.149899$ in this compendium's.

\emph{Original comparisons} (same tree and shift).
\begin{itemize}
\item \emph{One-sided.} The callable's down-duration is 2.142 and its up-duration 2.335.
\item \emph{Callable, putable and straight.} Effective durations are 2.239 for the callable, 2.816 for the putable and 3.517 for the straight bond. Convexities (compendium convention) are $-77.15$, $+11.13$ and $+16.36$. With $C=-77.15$, $-D\Delta r+\tfrac12C\Delta r^2$ returns $V_\pm$ exactly at the 25 basis point shift; for the straight bond at $\pm50$ basis points it is within 0.02 basis points of full revaluation, while putting the half-scaled figure into the same formula misses by about 1 basis point. Both options shorten duration, because each makes an early redemption at par likely in some states. The call adds negative convexity; the put keeps convexity positive.
\item \emph{Key rates.} The straight bond's key-rate durations, 0.010, 0.020, 0.031 and 3.455, sum to its effective duration (3.517 within 0.005); almost all its exposure sits at the four-year rate. The callable's, 0.443, 0.143, 0.166 and 1.539, move exposure to shorter rates, because the bond may be redeemed early. They sum to 2.290 rather than 2.239: a single-rate shift changes exercise decisions, so key-rate durations of an option-embedded bond are not exactly additive at a 25 basis point shift.
\end{itemize}

\qffield{Conditions} Measures depend on the rate model, the volatility, the shift size and the constant-OAS assumption; near an exercise boundary they change with the size of the shift (\hyperref[meth:curve-based-empirical-duration]{curve-based measures} shows a zero-volatility example). Key-rate durations depend on the chosen vertices.

\qffield{Pitfalls} Holding the price, rather than the OAS, constant when shifting the curve. Hedging a callable with its effective duration alone through a large move. Summing key-rate durations of an option-embedded bond and expecting the parallel-shift figure exactly.

\qffield{See also} \hyperref[meth:yield-duration-pvbp]{yield duration}; \hyperref[meth:yield-curve-factor-exposures]{factor exposures}.

\end{qfmethod}

\begin{qfmethod}[Capped and floored floating-rate bonds: valuation on a rate tree with coupons set in advance and paid in arrears, and the cap and floor as separate options]\label{meth:capped-floored-floater-valuation}
\qffield{Statement} A floater's coupon is set from the reference rate at the start of each period and paid at its end, in arrears. On the tree, the value at a node is $\bigl[\tfrac12(V_H+V_L)+100\,c(r^*)\bigr]/(1+r^*)$, where $c(r^*)$ is the coupon set from the node rate $r^*$ \citep[chap.~9, sec.~IX, pp.~244--247]{fixed-income-analysis-2007}. A floater paying the reference rate flat is worth par at every node. With a cap $K$, $c=\min(r^*,K)$; with a floor $F$, $c=\max(r^*,F)$. ``The higher the cap, the closer the capped floater will trade to par'' (p.~246). The lattice can value floaters with option-type features, caps, floors and swaptions in the same way \citep[chap.~36, p.~832]{s7-fabozzifixed}.

\emph{As options.} A capped floater is a floater less a cap that the investor has sold, a strip of caplets paying $\max(r-K,0)$ in arrears. A floored floater is a floater plus a floor, paying $\max(F-r,0)$. A collared floater combines both (Interpretation; cap and floor contracts in \hyperref[meth:black-model-rates-currencies]{Black's model}).

\qffield{Practice} \emph{Reproduced} (Fixed Income Analysis, Exhibits~21--22, the Handbook's tree): a floater at the one-year rate is worth 100; with a 7.25 percent cap 99.724; with a 7.75 percent cap 99.858.

\emph{Original} (same tree). The 7.25 percent cap valued alone is worth 0.276, so $100-0.276=99.724$. A 4.5 percent floor raises the floater to 100.992; the floor valued alone is worth 0.992. A collared floater with the 4.5 percent floor and the 7.25 percent cap is worth $100+0.992-0.276=100.716$.

\qffield{Conditions} The reference rate and the discount rate are the same tree rate, and the floater pays no margin; with a quoted margin or a different discount curve the uncapped floater is not at par (\hyperref[meth:floating-money-market-yields]{floating-rate measures}).

\qffield{Pitfalls} Applying the cap to the rate observed at the payment date instead of the reset date. Discounting the period's coupon at the next node's rate. Valuing the cap as if the investor held it.

\qffield{See also} \hyperref[meth:callable-putable-bond-valuation]{callable bonds}; \hyperref[meth:option-adjusted-spread]{OAS}.

\end{qfmethod}
